% Additions to Section 2 of the master: reflective seams, half-Hopf pieces
% and tubes in a conformally flat exterior.
% Drafted for Edward B. Baker III by Claude (Anthropic, model claude-fable-5-2,
% Cowork session, reasoning effort max as configured), 26 November 2026, from
% papers/master/brainstorm/notes/claude/PATHS_FORWARD_AFTER_ATTAINMENT_AND_ENTRY_TRIALS_20261126.md
% at repository state de2125d. The formulas were tested by the scripts filed
% with that note. No independent reading.

\paragraph{The two sides of a reflective seam.}
The seam test gives a sign. A second necessary condition measures the
energy on each side. Let \((Q,[g])\) be a closed oriented conformal
four-manifold and \(V\subset Q\) a compact domain with smooth boundary.
Call \(\partial V\) \emph{reflective} if \([g]\) is invariant under an
involution of a neighborhood of \(\partial V\) which fixes \(\partial V\)
pointwise and exchanges its two sides. As for the identity doubles
above, \(\partial V\) is then totally geodesic in an invariant
representative, hence umbilic. Conversely an umbilic hypersurface is
reflective wherever the class is locally conformally flat near it: in
flat charts it lies in round spheres, and the inversions in these
spheres agree on overlaps. For reflective \(\partial V\) the double
\(\mathcal D(V)=V\cup_{\partial V}\overline V\) carries the smooth
reflected class. Put
\[
 e_V=\tfrac12\big(E_+(V)+E_-(V)\big).
\]
If a reflective seam separates the homotopy sphere \(M\) into \(V\) and
\(V'\), then \(e=e_V+e_{V'}\). In a representative \(h\) for which an
umbilic \(\partial V\) is totally geodesic, Chern--Gauss--Bonnet on
\(V\) has no boundary term, so that
\[
 e_V=T\chi(V)+\tfrac14\int_V|Z_h|^2\,dV_h-\tfrac1{48}\int_VR_h^2\,dV_h .
\]

\begin{proposition}[Doubling at a reflective seam]\label{prop:doubling}
Let \(\partial V\) be reflective and let \(Y_{\mathcal D}\) be the Yamabe
constant of the reflected class on \(\mathcal D(V)\). Then
\[
 e_V\ \ge\ T\chi(V)-\frac{Y_{\mathcal D}^2}{96}.
\]
If \(Y_{\mathcal D}>0\) and \(b_1(\mathcal D(V))>0\), then
\(e_V\ge T\chi(V)\). If \(Y_{\mathcal D}>0\) and
\(b_2(\mathcal D(V))>0\), then \(e_V\ge\tfrac23T\chi(V)\).
\end{proposition}
\begin{proof}
The reflection of \(\mathcal D(V)\) reverses orientation and exchanges
the two Weyl blocks. Hence \(E_\pm(\mathcal D(V))=2e_V\) and
\(\tau(\mathcal D(V))=0\), so that \(b_2^+=b_2/2\). Since
\(\chi(\partial V)=0\), one has \(\chi(\mathcal D(V))=2\chi(V)\).
Take a volume-one Yamabe representative of the reflected class, of
constant scalar curvature \(Y_{\mathcal D}\)
\cite{Schoen1984Yamabe,LeeParker1987Yamabe}. Chern--Gauss--Bonnet reads
\[
 16\pi^2\chi(V)=4e_V+\frac{Y_{\mathcal D}^2}{24}
       -\frac12\int_{\mathcal D(V)}|Z|^2\,dV ,
\]
which gives the first inequality. For a class of positive Yamabe
constant on a closed oriented four-manifold, Gursky's inequalities are
\(E_+\ge2\pi^2(2\chi+3\tau)\) when \(b_1>0\) and
\(E_+\ge\tfrac{4\pi^2}{3}(2\chi+3\tau)\) when \(b_2^+>0\)
\cite{Gursky1998WeylFunctional}; see
\cite[Theorems~A and~C]{Fu2018FourManifolds} for the statements, in the
four-index norm. With \(E_+=2e_V\), \(\tau=0\) and
\(\chi=2\chi(V)\) they give the other two bounds.
\end{proof}

A \emph{half-Hopf piece} of length \(\ell\) in \((Q,[g])\) is a compact
domain conformal to \(S^1_\ell\times\overline{S^3_+}\), the product of a
circle of length \(\ell\) and a closed hemisphere of the unit
three-sphere, by a map which extends conformally to a neighborhood in
\(S^1_\ell\times S^3\). It carries no Weyl energy, and its boundary
\(S^1_\ell\times S^2\) is reflective.

\begin{corollary}[Half-Hopf pieces]\label{cor:half-hopf}
Let a class on the homotopy sphere \(M\) contain \(k\ge1\) pairwise
disjoint half-Hopf pieces, let \(V\) be the closure of their complement,
and let \(Y_{\mathcal D}\) be the Yamabe constant of \(\mathcal D(V)\).
Then
\[
 e\ \ge\ 2T-\frac{Y_{\mathcal D}^2}{96}.
\]
If \(Y_{\mathcal D}>0\), then \(e\ge\tfrac43T\). If
\(Y_{\mathcal D}>0\) and \(k\ge2\), then \(e\ge2T\). In particular an
entering class which contains two disjoint half-Hopf pieces has
\(Y_{\mathcal D}\le-\sqrt{96(2T-e)}\).
\end{corollary}
\begin{proof}
Here \(e_V=e\) and \(\chi(V)=2\). The double is connected. By
Mayer--Vietoris \(b_1(\mathcal D(V))\ge k-1\), and
\(b_2(\mathcal D(V))=\chi(\mathcal D(V))-2+2b_1(\mathcal D(V))\ge2\).
Apply Proposition~\ref{prop:doubling}. If \(Y_{\mathcal D}\ge0\) and
\(k\ge2\), its first two bounds give \(e\ge2T\), which proves the last
statement.
\end{proof}

\paragraph{Tubes in a conformally flat exterior.}
Half-Hopf pieces occur in a conformally flat exterior from which
\(\Sigma_P\) is obtained by circle surgeries. Let \(F\) be a compact
hyperbolic surface with nonempty geodesic boundary, and let
\(\Gamma\subset SO_0(2,1)\subset SO_0(5,1)\) be the Fuchsian group of
its completion. On \(S^4\) the limit set \(\Lambda\) of \(\Gamma\) lies
in a round circle whose complement is conformal to
\(\mathbb H^2\times S^2\). The quotient
\(N=(S^4\setminus\Lambda)/\Gamma\) is a flat conformal structure on the
boundary of \(F\times D^3\), a connected sum of \(n=1-\chi(F)\ge2\)
copies of \(S^1\times S^3\), with fundamental group \(\Gamma\). For
\(\gamma\ne1\) in \(\Gamma\), of translation length \(\ell\), let \(r\)
be the signed distance from its axis in \(\mathbb H^2\) and \(t\) the
arclength along the axis. Then
\begin{equation}\label{eq:hopf-coordinates}
 dt^2+d\beta^2+\cos^2\!\beta\,g_{S^2}
 =\cos^2\!\beta\,\big(dr^2+\cosh^2\!r\,dt^2+g_{S^2}\big),
 \qquad \cosh r=\sec\beta ,
\end{equation}
and the left side is the product metric of \(S^1_\ell\times S^3\) at
the point \(y=\sin\beta\,E+\cos\beta\,\omega\) of
\(S^3\subset\mathbb R\oplus\mathbb R^3\), \(\omega\in S^2\). So the
quotient of \(S^4\) minus the two fixed points by
\(\langle\gamma\rangle\) is the Hopf manifold \(S^1_\ell\times S^3\).
It contains \((S^4\setminus\Lambda)/\langle\gamma\rangle\), and the two
sides of the axis are the hemispheres \(\pm\langle y,E\rangle>0\).
The part of \(N\) over the closed funnel beyond a boundary geodesic of
\(F\), with its circle at infinity, is thus a half-Hopf piece.

\begin{proposition}[One tube, closed on the standard sphere]
\label{prop:hopf-closing}
Let \(c\subset N\) be an embedded closed curve freely homotopic to
\(\gamma\ne1\) in \(\Gamma\), of translation length \(\ell\), and let
\(\mathcal T\) be a closed tubular neighborhood of \(c\). Replace
\(\mathcal T\) by \(B\cong D^2\times S^2\), with a conformal class on
the surgered manifold which is the flat one outside \(B\).
\begin{enumerate}
\item \(E_+(B)=E_-(B)\). Write \(c_B\) for this number.
\item Suppose that the lift of \(\mathcal T\) to
\(S^1_\ell\times S^3\) lies in \(S^1_\ell\times S^3_+\) for an open
hemisphere \(S^3_+\). Then there is a reflection-invariant class on
\(S^2\times S^2\) with \(E_+=E_-=2c_B\), of Yamabe constant
\(Y_{\mathcal D}\), such that
\[
 c_B\ \ge\ 2T-\frac{Y_{\mathcal D}^2}{96},\qquad\text{and}\qquad
 c_B\ \ge\ \tfrac43T\quad\text{if }Y_{\mathcal D}>0 .
\]
\end{enumerate}
\end{proposition}
\begin{proof}
Since \(c\) is freely homotopic to \(\gamma\), the tube lifts
homeomorphically to a tube \(\widetilde{\mathcal T}\) in
\((S^4\setminus\Lambda)/\langle\gamma\rangle\) about a closed curve
which generates the fundamental group of the Hopf manifold. Homotopic
embedded circles in a four-manifold are isotopic, so this curve is
isotopic to the circle factor. Replace \(\widetilde{\mathcal T}\) by
\(B\) with the same class. Surgery on the circle factor of
\(S^1\times S^3\) gives \(S^4\) for either framing, because the Gluck
twist of \(S^1\times S^2\) extends over \(S^1\times D^3\). The result
is therefore a conformal class on \(S^4\) which is flat outside \(B\).
Its chiral energies are \(E_\pm(B)\), and they agree by the signature
formula. Under the hypothesis of (2) the complementary closed
hemisphere gives a half-Hopf piece
\(S^1_\ell\times(S^3\setminus S^3_+)\) of this class on \(S^4\). The
closure \(V\) of its complement is the complement of an open tubular
neighborhood of a circle in \(S^4\). Hence \(V\cong D^2\times S^2\)
and \(\mathcal D(V)\cong S^2\times S^2\), with
\(E_\pm(\mathcal D(V))=2c_B\). Corollary~\ref{cor:half-hopf} with
\(k=1\) gives the bounds.
\end{proof}

The product of two unit spheres has \(E_++E_-=16T/3\), and Kobayashi
conjectured that this is the infimum of the Weyl integral over all
metrics on \(S^2\times S^2\) \cite{Kobayashi1985Conformal}. The
conjecture holds among classes of positive Yamabe constant, by
Gursky's inequality, and among classes whose self-dual harmonic form
has no zero \cite[Theorem~B]{LeBrun2013DelPezzo}. It is open in
general, and the analogous bound fails for classes of negative Yamabe
constant on connected sums of many copies of \(S^2\times S^2\)
\cite[Theorem~A]{LeBrun2023Balance}. A filling with \(c_B<\tfrac43T\)
in Proposition~\ref{prop:hopf-closing}(2) has \(Y_{\mathcal D}<0\) and
gives \(E_++E_-=4c_B<16T/3\) on \(S^2\times S^2\). It would be a
counterexample.

\paragraph{Consequence for circle surgeries.}
Realize \(\Sigma_P\) from \(N\) by surgery on disjoint tubes
\(\mathcal T_1,\dots,\mathcal T_n\) about curves in the classes of the
relators, with fillings \(B_i\) and a class which is flat outside
them. Then \(e=\sum_i c_{B_i}\), and since \(n\ge2\) entry requires
\(c_{B_i}<T\) for some \(i\). Suppose that the lift of every tube lies
over an open hemisphere of its Hopf manifold. If each associated class
on \(S^2\times S^2\) satisfies Kobayashi's bound, then
Proposition~\ref{prop:hopf-closing} gives
\[
 e\ \ge\ \tfrac43\,nT\ >\ 2T .
\]
Entry by such a construction therefore requires a counterexample, and
indeed a class on \(S^2\times S^2\) with \(E_++E_-<4T\), three quarters
of the conjectured infimum. The hypothesis holds whenever a tube lies
in \(\mathbb H^2\times S^2\) modulo \(\Gamma\) and its \(S^2\)
coordinate stays in an open hemisphere, since
\(\langle y,\nu\rangle=\cos\beta\,\langle\omega,\nu\rangle\) for
\(\nu\in\mathbb R^3\). It does not involve the presentation, and the
conclusion applies equally to the trivial presentation drawn in this
way. If the tubes also avoid the half-Hopf pieces over the funnels,
these persist in the resulting homotopy sphere. When there are at
least two, an entering class has
\(Y_{\mathcal D}\le-\sqrt{96(2T-e)}\) on the double of their
complement by Corollary~\ref{cor:half-hopf}, while
\(Y([g])\ge\sqrt{48(2T-e)}\) on \(M\) by \eqref{eq:budgets}.

These bounds agree with surgery theory for the pointwise condition. By
Proposition~\ref{prop:entry} an entering class has a representative
with \(R-2\Ric(\nu,\nu)>0\) for every unit vector \(\nu\). This is
positivity of the \(1\)-curvature, which is preserved under surgeries
of codimension at least four \cite{Labbi1997Surgery}. A circle in a
four-manifold has codimension three.

\paragraph{Symmetric fillings and the retained waist.}
The hemisphere condition cannot be dropped. On a product
\(q+g_{S^2}\), with \(q\) a surface metric of curvature \(K\), both
\(W^\pm\) have eigenvalues \((K+1)(1/3,-1/6,-1/6)\)
\cite[The cost of a product circle cap]{Baker2026FixedDescent}. Thus
\begin{equation}\label{eq:product-energy}
 |W^\pm|_{\rm HS}^2=\frac{(K+1)^2}{6},\qquad
 E_\pm(U\times S^2)=\frac{2\pi}{3}\int_U(K+1)^2\,dA_q .
\end{equation}
Keep the part \(r\le r_0\) of the hyperbolic cylinder
\(dr^2+\cosh^2\!r\,dt^2\), \(t\in\mathbb R/\ell\mathbb Z\), and close
it by a disc \(\Delta\) whose metric matches in \(C^1\) along
\(r=r_0\). With the circle at \(r=-\infty\) this is a class on \(S^4\).
Since \((K+1)^2\ge4K\), the Gauss--Bonnet theorem on \(\Delta\) gives
\[
 E_\pm(\Delta\times S^2)\ \ge\ \frac{8\pi}{3}\int_\Delta K\,dA
 =\tfrac43T\Big(1+\frac{\ell}{2\pi}\sinh r_0\Big),
\]
with equality exactly when \(\Delta\) is a round cap of curvature one,
which requires \(\cosh2r_0=(2\pi/\ell)^2\) and is approached by smooth
fillings. For \(r_0\ge0\) the waist
\(r=0\) and its half-Hopf piece are retained and the cost is at least
\(\tfrac43T\), in agreement with Proposition~\ref{prop:hopf-closing}.
For \(r_0<0\) the filling swallows the waist and a collar of area
\(\ell\sinh|r_0|\), and the round cap costs
\(\tfrac43T\big(1-\sqrt{(1-\ell^2/4\pi^2)/2}\,\big)\). The constant
\(\tfrac43T\) thus belongs to a retained waist with its funnel. A
product filling of cost below \(\tfrac43T(1-f)\) has to swallow a
collar of area more than \(2\pi f\), and \(2\pi\) is the area of a pair
of pants.
